% templates/resume.tex
% Canonical resume template. Layout spec in templates/RESUME_FORMAT.md.
%
% Build:   pdflatex resume.tex      (or: latexmk -pdf resume.tex)
% Output:  one-page PDF, US Letter, 0.5in margins, Helvetica (Arial match).
%
% ---------------------------------------------------------------------------
% FIRST RUN: this file ships with {{PLACEHOLDERS}}. The `setup` skill fills the
% header and Education block from your answers, once. After that this file is
% YOUR canonical resume and the placeholders are gone.
%
% It compiles as-shipped, placeholders and all, so you can run
%   pdflatex templates/resume.tex
% right after cloning to confirm your LaTeX install works before entering
% anything real. (Placeholder names avoid underscores for exactly this reason:
% `_` is a math-mode character and would break the build.)
%
% Per-job resumes are produced by cloning this file into
% outputs/<company-role>/<LastName>-Resume.tex and swapping only the Relevant
% Experience and Work Experience entries (and their ordering) to suit the
% posting. Do not redesign the layout; change content only.
%
% If you already have a resume you like, tell the setup skill and it will
% reshape this template to match your layout before filling anything in.
% ---------------------------------------------------------------------------

\documentclass[11pt,letterpaper]{article}

\usepackage[T1]{fontenc}
\usepackage[utf8]{inputenc}
\usepackage[margin=0.5in]{geometry}
\usepackage{helvet}                 % Helvetica: the closest metric match to Arial
\renewcommand{\familydefault}{\sfdefault}
\usepackage[shortlabels]{enumitem}
\usepackage{hyperref}
\hypersetup{
  colorlinks=true,
  urlcolor=blue,
  linkcolor=blue,
  pdfauthor={{{FULL-NAME}}},
  pdftitle={{{FULL-NAME}} - Resume}
}

\pagestyle{empty}
\setlength{\parindent}{0pt}

% Tight bullet lists: 11pt round bullets, minimal vertical air.
\setlist[itemize]{leftmargin=1.5em, topsep=2pt, itemsep=1.5pt, parsep=0pt,
                  label=$\bullet$}

% --- Helpers -----------------------------------------------------------------
% Section header: 12pt bold with a full-width bottom rule.
\newcommand{\sectionheader}[1]{%
  \vspace{6pt}%
  {\fontsize{12}{14}\selectfont\bfseries #1}\par
  \vspace{1pt}\hrule height 0.6pt\vspace{4pt}%
}

% Entry: bold title line, then org (left) with date/term right-aligned.
% Usage: \entry{Title}{Organization line}{Date or term}
\newcommand{\entry}[3]{%
  \vspace{5pt}%
  \textbf{#1}\\
  #2 \hfill \textit{#3}%
}

\begin{document}

% --- Name + contact ----------------------------------------------------------
% Contact line format: phone | email | portfolio.
% No portfolio site? Delete the third segment AND the \quad\textbar\quad before
% it. LinkedIn can take its place if you would rather link that.
\begin{center}
  {\fontsize{18}{22}\selectfont\bfseries {{FULL-NAME}}}\\[4pt]
  {{PHONE}} \quad\textbar\quad \href{mailto:{{EMAIL}}}{{{EMAIL}}} \quad\textbar\quad \href{https://{{PORTFOLIO}}}{{{PORTFOLIO}}}
\end{center}

% --- Education ---------------------------------------------------------------
% Drop any line that does not apply to you. Nothing here is mandatory except the
% school, degree, and graduation term.
\sectionheader{Education}
\textbf{{{SCHOOL}}} \hfill \textit{{{GRADUATION-TERM}}}\\
\textit{{{DEGREE}}}\\
{{HONORS-LINE}}\\
{{CREDENTIALS-LINE}}\\
\textbf{Activities:} {{ACTIVITIES}}\\
\textbf{Technical Skills:} {{TECHNICAL-SKILLS}}

% --- Relevant Experience -----------------------------------------------------
% Everything below this line is replaced per job by the build-application skill.
% The two entries here exist only to show the pattern. Delete them once you have
% run your first build; they are not your content.
\sectionheader{Relevant Experience}

\entry{{{JOB-TITLE}} - {{TEAM-OR-LAB}}}%
{{{ORGANIZATION}}}{{{DATE-RANGE}}}
\begin{itemize}
  \item Bullet drawn from a Resume Bullet Candidate in the matching accomplishment doc.
  \item Second bullet. Strong verb first, technical detail, impact where a real number exists.
  \item Third bullet. Two or three per entry reads best at this density.
\end{itemize}

% --- Work Experience ---------------------------------------------------------
\sectionheader{Work Experience}

\entry{{{JOB-TITLE}} - {{EMPLOYER}}}%
{{{LOCATION}}}{{{DATE-RANGE}}}
\begin{itemize}
  \item Employment that shows reliability and working with people belongs here.
  \item Engineering-relevant work belongs in Relevant Experience above.
\end{itemize}

\end{document}
